\dish{味噌ラーメン}
\hiragana{みそラーメン}
\altdish{Miso ramen}
\altdish{Ramen, miso (Sapporo style)}
\serves{4}
\prep{broth 2\fracH~h (mostly unattended) + 30~min}
\source{openai-chatgpt}


\stage{Broth}

\begin{ingredients}
  \ingrS{1}{kg}{chicken wings, backs or carcasses}
  \ingr{2}{L}{water}
  \ingr{20}{g}{ginger, sliced}
  \ingr{3}{cloves}{garlic, crushed}
  \ingr{1}{}{small onion, halved}
  \ingr{2}{}{spring onions}
  \ingrS{4}{}{dried shiitake}
  \ingr{15}{g}{kombu}
  \ingr{15}{g}{katsuobushi (optional)}
\end{ingredients}

\begin{recipe}
  \begin{enumerate}

  \item Cover the chicken with cold water, boil 3~min. Drain, rinse
    chicken and pot.

  \item Return chicken to the pot with the water, ginger, garlic,
    onion, spring onions and shiitake. Simmer gently, partially
    covered, 2--2\fracH~h. Top up with water as needed.

  \item Add kombu for the last 20--30~min, keeping below a simmer
    (80--85~\textdegree{}C). Remove kombu.

  \item Off heat, steep katsuobushi 3--5~min. Strain. Aim for
    1.4--1.6~L; top up with hot water. Don't salt.

  \end{enumerate}
\end{recipe}


\stage{Miso tare}

\begin{ingredients}
  \ingr{80}{g}{red miso}
  \ingr{60}{g}{white miso}
  \ingr{2}{\tbsp}{Japanese soy sauce}
  \ingr{2}{\tbsp}{mirin}
  \ingr{1}{\tbsp}{sake}
  \ingr{1}{\tbsp}{toasted sesame seeds, ground}
  \ingr{1}{\tsp}{sugar}
  \ingr{1}{\tsp}{doubanjiang (optional)}
  \ingr{1}{\tbsp}{nerigoma or tahini (optional)}
\end{ingredients}

\begin{recipe}
  \begin{enumerate}

  \item Mix to a smooth paste.

  \end{enumerate}
\end{recipe}


\stage{Soup}

\begin{ingredients}
  \ingr{2}{\tbsp}{neutral oil}
  \ingrS{250}{g}{minced pork, not too lean}
  \ingr{4}{cloves}{garlic, minced}
  \ingr{25}{g}{ginger, grated}
  \ingrS{1}{}{onion, thinly sliced}
  \ingr{150}{g}{cabbage, thinly sliced}
  \ingrS{150}{g}{bean sprouts}
  \ingr{1}{\tbsp}{toasted sesame oil}
\end{ingredients}

\begin{recipe}
  \begin{enumerate}

  \item In a wok or large pan, brown the pork in the oil over
    medium-high heat, 4--5~min.

  \item Add garlic, ginger and onion, 2--3~min.

  \item Lower to medium. Add the tare and stir-fry 1--2~min until
    fragrant and slightly darker, without burning.

  \item Add cabbage and bean sprouts, 2~min; keep some crunch.

  \item Gradually add the hot broth, stirring to dissolve the miso.
    Simmer 3--5~min. Adjust with soy sauce or water. Stir in the
    sesame oil.

  \end{enumerate}
\end{recipe}


\stage{Assembly}

\begin{ingredients}
  \ingr{4}{}{eggs}
  \ingr{500--600}{g}{\hyperref[ラーメン]{fresh ramen noodles} (or 320--400~g dried)}
  \ingr{100}{g}{sweetcorn}
  \ingr{4}{}{spring onions, finely sliced}
  \ingr{20}{g}{butter (optional)}
  \ingr{}{}{nori, chilli oil (optional)}
\end{ingredients}

\begin{recipe}
  \begin{enumerate}

  \item Lower cold eggs into boiling water, 6\fracH--7~min. Ice water,
    peel, halve.

  \item Cook the noodles; divide among 4 hot bowls. Ladle over the
    soup, sharing out pork and vegetables. Top with egg, corn, spring
    onion, a pat of butter, and nori or chilli oil.

  \end{enumerate}
\end{recipe}

\notes Broth and tare keep several days refrigerated; dinner then
takes about 30~min. A single robust awase miso can replace the red
and white.

%\accord{}
